\documentclass[11pt]{article}
\usepackage[margin=1in]{geometry}
\usepackage{booktabs}
\usepackage{hyperref}
\usepackage{xcolor}
\usepackage{microtype}

\title{Subtraction Is Conditional:\\
An Interaction-Focused Replication of Deletion-First Refactoring Prompts}
\author{The Subtraction Research Project}
\date{August 1, 2026}

\begin{document}
\maketitle

\begin{abstract}
Agentic coding systems often produce additive patches even when a task is framed
as a refactor. We tested whether a previously observed deletion-first component
effect would replicate across two larger deterministic refactor fixtures. The
registered design crossed two tasks, three model/effort pairs, eight combinations
of task-type, delete-first, and semantic net-line-budget components, and five
repetitions: 240 planned cells. All 240 selected completed receipts passed the
behavior oracle. The deletion-first factor was negative on the shared-normalizer
fixture for GPT-5.6 Luna and Grok 4.5, but null for Composer; it was null or
positive on the dead-compatibility-path fixture for all three models. The result
is therefore task- and model-conditional subtraction, not a universal standalone
D component. Six initial adapter failures are retained as separate transport
provenance and excluded from completed comparisons. No priced cost claim is made.
\end{abstract}

\section{Question and boundary}
This study asks whether prompt components can shift the shape of an agentic
refactor toward subtraction without sacrificing behavior. The claim tested here
is deliberately narrow: does the deletion-first component produce a stable
raw-net effect under a larger-refactor replication? A negative raw-net value
means fewer net lines in the resulting patch; it is not, by itself, evidence of
better engineering or lower cost.

\section{Registered design}
The wave used two deterministic tasks:
\begin{itemize}
\item \texttt{refactor-shared-normalizer};
\item \texttt{refactor-dead-compatibility-path}.
\end{itemize}
Each was crossed with the full eight-arm T/D/B factorial, three frozen
model/effort pairs, and repetitions 1--5. Thus the planned size was
\[
  2 \times 3 \times 8 \times 5 = 240.
\]
The models were GPT-5.6 Luna (maximum), Grok 4.5 (high), and Composer 2.5
(default). Completed comparisons require matched repetition IDs, completed
adapter status, and passing behavior oracles. Initial failures remain visible
in provenance artifacts rather than being silently treated as missing success.

\section{Results}
Table~\ref{tab:effects} reports the factorial D main effect and the matched
D-only treatment-minus-neutral contrast, in raw-net lines. These are descriptive
means over the five repetitions.

\begin{table}[ht]
\centering
\caption{Deletion-first effects by task and model. Negative values are more subtractive.}
\label{tab:effects}
\begin{tabular}{llrr}
\toprule
Task & Model & D factorial & D-only contrast \\
\midrule
Shared normalizer & Composer 2.5 & -0.65 & 0.0 \\
Shared normalizer & GPT-5.6 Luna & -3.10 & -6.4 \\
Shared normalizer & Grok 4.5 & -3.25 & -10.0 \\
Dead compatibility path & Composer 2.5 & 0.0 & 0.0 \\
Dead compatibility path & GPT-5.6 Luna & +0.30 & +0.8 \\
Dead compatibility path & Grok 4.5 & +2.0 & +3.4 \\
\bottomrule
\end{tabular}
\end{table}

All selected observations passed their behavior checks: 240 of 240. The
negative effect replicated across two models on one fixture, but not across
both fixtures. This pattern is consistent with an interaction between task
structure and model behavior. It does not establish a portable one-component
mechanism, statistical significance, causality, or universality across
repositories and maintenance tasks.

\section{Provenance and cost boundary}
The live run recorded six initial adapter failures. They were retried from
fresh isolated roots and remain separate provenance; the completed dataset
contains 240 receipts and excludes those six non-completed attempts from
completed-run statistics. Token fields are preserved as descriptive telemetry.
The planned priced multi-turn cost wave was deferred because the patch-shape
effect did not survive both larger fixtures. No dollar-saving claim follows.

\section{Conclusion}
Deletion-first/reference-proof prompting can produce materially more
subtractive patches on some larger refactor structures. In this replication,
the effect is conditional on task and model. The next evidence gate is an
independently specified refactor replication on a new task or repository before
making broader claims or modeling priced trajectories.

\section*{Reproducibility}
The protocol, findings, frozen manifest, receipts, factor effects, summary,
and wave metadata are committed in the repository:
\begin{itemize}
\item \url{../research/phase-2/LARGER_REFACTOR_PROTOCOL.md}
\item \url{../research/phase-2/LARGER_REFACTOR_FINDINGS.md}
\item \url{../research/phase-2/design-larger-refactor-screen-r5.json}
\item \url{../research/phase-2/live/larger-refactor-r5-results.json}
\end{itemize}
Build this paper from the repository root with Tectonic using the
\texttt{paper/subtraction-study.tex} source, or from the \texttt{paper/}
directory with \texttt{latexmk -pdf subtraction-study.tex}.

\end{document}
